\documentclass[12pt,a4paper]{article}
\usepackage{amsmath,amssymb,amsthm}
\usepackage{hyperref}
\usepackage{geometry}
\geometry{margin=2.5cm}
\usepackage{enumitem}
\usepackage{booktabs}

\newtheorem{theorem}{Theorem}[section]
\newtheorem{lemma}[theorem]{Lemma}
\newtheorem{corollary}[theorem]{Corollary}
\newtheorem{proposition}[theorem]{Proposition}
\theoremstyle{definition}
\newtheorem{definition}[theorem]{Definition}
\theoremstyle{remark}
\newtheorem{remark}[theorem]{Remark}

\title{The Fine Structure Constant from Recognition Science:\\
$\alpha^{-1} \approx 137.035$ from $D = 3$ Cube Geometry}

\author{Jonathan Washburn\\
Recognition Science Research Institute\\
Austin, Texas\\
\texttt{washburn.jonathan@gmail.com}}

\date{March 2026}

\begin{document}
\maketitle

\begin{abstract}
We derive the fine structure constant $\alpha$ from the Recognition
Science (RS) framework.  The derivation is additive, proceeding in
three steps from $D = 3$ cube geometry:
(i)~a geometric seed $\alpha_{\mathrm{seed}} = 4\pi \times 11
\approx 138.230$, where $11 = E_{\mathrm{passive}}$ is the number
of passive edges of the $D = 3$ cube;
(ii)~a gap subtraction $f_{\mathrm{gap}} = w_8 \cdot \ln\varphi
\approx 1.199$, where $w_8$ is the parameter-free 8-tick gap
weight (a closed-form algebraic number);
(iii)~a curvature correction $\delta_\kappa = -103/(102\,\pi^5)
\approx -0.0033$, where $102 = 6 \times 17$ (cube faces times
wallpaper groups) and $103 = 102 + 1$ (Euler closure).  The
result:
$\alpha^{-1} = \alpha_{\mathrm{seed}} - f_{\mathrm{gap}} -
\delta_\kappa \approx 137.035$.
The RS prediction is for the Thomson limit ($Q^2 = 0$), giving
$\alpha^{-1} \approx 137.0349$, which agrees with the CODATA 2024
value $137.035\,999\,206(11)$ at the $8.1\;\mathrm{ppm}$ level.
All structural results, including the closed-form expression for
$w_8$, are derived from the RS axioms~\cite{RS_Axioms}.
\end{abstract}

%% ============================================================
\section{Introduction}
\label{sec:intro}
%% ============================================================

The fine structure constant
$\alpha = e^2/(4\pi\epsilon_0\hbar c) \approx 1/137.036$ governs
all electromagnetic interactions.  It is one of the most precisely
measured constants in physics~\cite{Morel2020}:
\begin{equation}
  \alpha^{-1} = 137.035\,999\,206 \pm 0.000\,000\,011
  \quad (\text{CODATA 2024}).
\end{equation}
Feynman called it ``one of the greatest damn mysteries of
physics''---a pure number that no theory had explained.

In the Standard Model, $\alpha$ is a free parameter.  RS~\cite{RS_Axioms} derives
it from the geometry of the $D = 3$ cube, with zero free
parameters.

%% ============================================================
\section{Cost Functional Foundation}
\label{sec:cost}
%% ============================================================

RS derives all physical law from three axioms~\cite{RS_Axioms}:
\begin{enumerate}[label=(A\arabic*),nosep]
  \item $J(1) = 0$ \hfill\textit{(Normalization)}
  \item $J(xy) + J(x/y) = 2J(x)J(y) + 2J(x) + 2J(y)$ for all $x,y>0$
  \hfill\textit{(Recognition Composition Law)}
  \item $J''_{\log}(0) = 1$, where $J_{\log}(t) := J(e^t)$
  \hfill\textit{(Calibration)}
\end{enumerate}

\begin{theorem}[Cost Uniqueness~\cite{RS_Axioms}]
\label{thm:J}
The unique continuous solution is $J(x) = \tfrac{1}{2}(x+x^{-1})-1$,
satisfying $J(x) \geq 0$ with equality iff $x = 1$.
\end{theorem}

\begin{proof}[Proof sketch]
Setting $H(t) = 1 + J(e^t)$ gives the d'Alembert equation
$H(s+t)+H(s-t)=2H(s)H(t)$ with $H(0)=1$ and $H''(0)=1$.
The Acz\'el classification~\cite{Aczel1966} forces
$H(t)=\cosh t$, giving $J(x) = \tfrac{1}{2}(x+x^{-1})-1$.
\end{proof}

The golden ratio $\varphi = (1+\sqrt{5})/2$ is the unique positive root of
$x^2 = x + 1$, forced by ledger self-similarity.  The fine-structure constant
enters through the cost of electromagnetic recognition events: each photon
exchange traverses the 8-tick epoch, incurring a $J$-cost that reduces the
effective coupling from the geometric seed.

%% ============================================================
\section{Cube Geometry Ingredients}
\label{sec:ingredients}
%% ============================================================

The derivation uses three quantities, all determined by $D = 3$.

\begin{definition}[Geometric Seed]
\label{def:seed}
\begin{equation}
  \alpha_{\mathrm{seed}} = 4\pi \times E_{\mathrm{passive}}
  = 4\pi \times 11 = 44\pi \approx 138.230,
\end{equation}
where $E_{\mathrm{passive}} = 2D(D-1) - 1 = 12 - 1 = 11$ is the
number of passive edges of the $D = 3$ cube.
\end{definition}

\begin{remark}
The $D = 3$ cube has $E = 2D(D-1) = 12$ edges.  One edge is
distinguished as the recognition axis (the direction along which
comparison occurs), leaving $11$ passive edges.  The factor $4\pi$
is the solid angle of the full sphere, representing isotropic
averaging over all orientations.
\end{remark}

\begin{definition}[Gap Weight]
\label{def:gap}
The 8-tick gap weight $w_8$ is the normalized DFT projection of
the $\varphi$-pattern onto the 8-tick epoch structure:
\begin{equation}
  w_8 = \frac{246 + 145\sqrt{2} - 102\sqrt{5} - 65\sqrt{10}}{7}
  \approx 2.4906.
\end{equation}
This is a parameter-free closed-form algebraic number derived from
the 8-tick epoch.  The gap subtraction is:
\begin{equation}
  f_{\mathrm{gap}} = w_8 \cdot \ln\varphi \approx 1.199.
\end{equation}
\end{definition}

\begin{remark}
The closed form for $w_8$ is a parameter-free algebraic number.
The derivation chain uses the Parseval identity for the $\varphi$-pattern DFT,
the phiDFT denominator formula, and a polynomial identity linking
the projection to the algebraic expression above.
\end{remark}

\begin{definition}[Curvature Correction]
\label{def:curvature}
\begin{equation}
  \delta_\kappa = -\frac{103}{102\,\pi^5} \approx -0.0033,
\end{equation}
where:
\begin{itemize}[nosep]
  \item $102 = 6 \times 17$: cube faces ($2D = 6$) times wallpaper
  groups ($17$, the number of distinct 2D crystallographic
  symmetries on each face).
  \item $103 = 102 + 1$: the Euler closure term---the extra $+1$
  accounts for the global topological constraint (the recognition
  seam that closes the cube into a torus).
  \item $\pi^5$: arises from the 5-dimensional phase-space
  integration over the compact directions ($2^D - D = 5$
  non-spatial modes of the 8-tick epoch).
\end{itemize}
\end{definition}

%% ============================================================
\section{The Main Result}
\label{sec:derivation}
%% ============================================================

\begin{proposition}[Inverse Fine Structure Constant --- Three-Term RS Formula]
\label{thm:alpha}
\begin{equation}
  \boxed{\alpha^{-1} = 4\pi \times 11 \;-\;
  w_8\,\ln\varphi \;-\; \delta_\kappa
  \;\approx\; 137.035.}
\end{equation}
\end{proposition}

\begin{proof}
Each term has a distinct physical origin:
\begin{enumerate}[label=(\roman*),nosep]
  \item \textbf{Seed} ($44\pi \approx 138.230$): The number of
  passive photon-exchange channels, weighted by the solid angle.
  This sets the ``bare'' inverse coupling at the recognition scale.

  \item \textbf{Gap subtraction} ($-w_8\,\ln\varphi \approx
  -1.199$): The $J$-cost penalty for an electromagnetic
  recognition event within one 8-tick epoch.  Each photon exchange
  must traverse the epoch, incurring a $J$-cost proportional to
  $\ln\varphi$ (the minimum ledger bit-cost) weighted by the
  8-tick projection $w_8$.

  \item \textbf{Curvature correction}
  ($-\delta_\kappa \approx +0.003$): The topological correction
  from closing the cube faces into a compact space.
\end{enumerate}

Numerically, using $w_8 = 2.490569\ldots$, $\ln\varphi = 0.481212\ldots$,
and $\pi^5 = 306.0197\ldots$:
\begin{align}
  44\pi &= 138.23008\ldots \\
  w_8\ln\varphi &= 1.19849\ldots \\
  \delta_\kappa &= -\frac{103}{102 \times 306.020} = -0.003300\ldots \\[4pt]
  \alpha^{-1} &= 138.23008 - 1.19849 - (-0.00330)
  = 137.03489\ldots \approx 137.0349.
\end{align}
The predicted value $\alpha^{-1} \approx 137.0349$ (Thomson limit) agrees
with the CODATA 2024 measurement $137.035\,999\,206(11)$ at the
$8.1\;\mathrm{ppm}$ level.  (See Remark~\ref{rem:precision} for the
significance of this agreement.)
\end{proof}

\begin{remark}[Thomson Limit and Precision Status]
\label{rem:precision}
The RS prediction is for the Thomson limit (zero momentum
transfer, $Q^2 = 0$), corresponding to the IR baseline of the
electromagnetic coupling.  The predicted value $\alpha^{-1} \approx
137.0349$ vs.\ the CODATA 2024 value $137.035\,999\,206(11)$
constitutes a fractional discrepancy of $8.1\;\mathrm{ppm}$.

\emph{Important caveat}: The experimental measurement precision is
$\sim 0.08\;\mathrm{ppb}$ ($8 \times 10^{-8}$), making the 8.1 ppm
discrepancy approximately $100{,}000 \times$ the measurement
uncertainty.  This means the three-term formula, while a remarkable
geometric approximation, is not in precise agreement with experiment.
The ``higher-order corrections'' that would close the gap are
currently unspecified within the RS framework.  The curvature
correction $\delta_\kappa = -103/(102\pi^5)$ in particular, while
motivated by the cube topology, has an ad hoc character (specifically
the ``Euler closure'' term $103 = 102 + 1$) that requires a more
rigorous derivation.  The three-term formula should be understood as
the leading approximation to $\alpha^{-1}$ from the RS geometric
structure, with the residual 8.1 ppm as an identified open problem.
\end{remark}

\begin{remark}[Exponential Resummation]
\label{rem:resum}
The additive three-term formula is the leading-order Taylor
expansion of the natural exponential form:
\begin{equation}
  \alpha^{-1}_{\mathrm{exp}} = \alpha_{\mathrm{seed}}
  \times \exp\!\Bigl(-\frac{w_8\ln\varphi}{\alpha_{\mathrm{seed}}}\Bigr)
  = 44\pi\,\exp\!\Bigl(-\frac{f_{\mathrm{gap}}}{44\pi}\Bigr)
  \approx 137.037,
\end{equation}
since $\exp(-x) \approx 1 - x$ for $x = f_{\mathrm{gap}}/\alpha_{\mathrm{seed}}
\approx 0.0087 \ll 1$.  The exponential form has the advantage of automatically
satisfying $\alpha^{-1} > 0$ for all values of $f_{\mathrm{gap}}$, and it achieves
$5.6\;\mathrm{ppm}$ agreement with CODATA \emph{without} the curvature correction
$\delta_\kappa$.  Whether the curvature correction should be applied additively
or absorbed into the exponential is an open question within RS.
\end{remark}

%% ============================================================
\section{The RS-Native Coupling}
\label{sec:native}
%% ============================================================

\begin{remark}[RS-Native Self-Similarity Coupling]
\label{thm:lock}
The RS-native self-similarity coupling (distinct from the
electromagnetic fine-structure constant) is:
\begin{equation}
  \alpha_{\mathrm{lock}} = \frac{1 - 1/\varphi}{2}
  = \frac{1}{2\varphi^2} \approx 0.191.
\end{equation}
This follows from $\varphi - 1 = 1/\varphi$:
$(1 - 1/\varphi)/2 = (2 - \varphi)/2 = 1/(2\varphi^2)$.
\end{remark}

\begin{remark}
$\alpha_{\mathrm{lock}} \approx 0.191$ is the tidal coefficient
of the ledger substrate.  It is the self-similarity exponent of
the $\varphi$-ladder, not the electromagnetic coupling constant.
The physical fine structure constant
$\alpha \approx 1/137 \approx 0.0073$ is derived from $D = 3$
cube geometry (Proposition~\ref{thm:alpha}).  The two quantities
differ by a factor of $\sim 26$ by design.
\end{remark}

%% ============================================================
\section{Origin Table}
\label{sec:origin}
%% ============================================================

All inputs trace to $D = 3$:

\begin{center}
\renewcommand{\arraystretch}{1.3}
\begin{tabular}{@{}lll@{}}
\toprule
\textbf{Factor} & \textbf{Value} &
\textbf{$D = 3$ origin} \\
\midrule
Passive edges & 11 & $2D(D-1) - 1 = 11$ \\
Solid angle & $4\pi$ & Isotropy of $S^2$ \\
Gap weight $w_8$ & 2.4906 & 8-tick DFT projection \\
$\ln\varphi$ & 0.4812 & Minimum $J$-cost bit \\
Wallpaper groups & 17 & 2D crystallography on faces \\
Cube faces & 6 & $2D = 6$ \\
Euler closure & $+1$ & Global topological seam (ad hoc) \\
Phase modes & 5 & $2^D - D = 5$ \\
\bottomrule
\end{tabular}
\end{center}

%% ============================================================
\section{Properties}
\label{sec:properties}
%% ============================================================

\begin{remark}[Bounds on $\alpha$]
\label{thm:bounds}
$\alpha > 0$, $\alpha < 1$, and
$\alpha^{-1} \in (137, 139)$.
To see this: $\alpha^{-1} = 44\pi - w_8\ln\varphi - \delta_\kappa$.
For the upper bound: $44\pi \approx 138.23$, $w_8\ln\varphi > 0$,
and $|\delta_\kappa| < 0.01$, so $\alpha^{-1} < 44\pi < 139$.
For the lower bound: $w_8 < 2.5$ and $\ln\varphi < 0.49$, so
$w_8\ln\varphi < 1.225$; also $|\delta_\kappa| < 0.004$; hence
$\alpha^{-1} > 44\pi - 1.225 + (-0.004) > 138.23 - 1.23 = 137.00$.
Therefore $\alpha^{-1} \in (137, 139)$ and $\alpha \in (1/139, 1/137)
\subset (0, 1)$.
\end{remark}

\begin{remark}[Running of $\alpha$]
\label{thm:running}
The fine structure constant runs with momentum transfer $Q^2$
according to the standard QED formula:
\begin{equation}
  \alpha(Q^2) = \frac{\alpha(0)}
  {1 - \Delta\alpha(Q^2)},
\end{equation}
where $\Delta\alpha$ includes leptonic and hadronic vacuum
polarization contributions.  At $Q^2 = M_Z^2$:
$\alpha^{-1}(M_Z) \approx 128$.  This running is a standard QED
result; the RS content is providing the input $\alpha(0)$ from
the three-term geometric formula.
\end{remark}

\begin{remark}
The RS derivation gives the Thomson-limit value.  The running
from $\alpha^{-1} \approx 137$ to $\alpha^{-1} \approx 128$ at
$M_Z$ is determined by the standard QED vacuum polarization, which
is itself derivable from the RS framework (charged lepton masses
from the $\varphi$-ladder).
\end{remark}

%% ============================================================
\section{Comparison}
\label{sec:comparison}
%% ============================================================

\begin{center}
\renewcommand{\arraystretch}{1.3}
\begin{tabular}{@{}lccc@{}}
\toprule
\textbf{Approach} & \textbf{Formula type} &
\textbf{$\alpha^{-1}$} &
\textbf{Free params} \\
\midrule
Eddington (1929) & Algebraic & 240 (wrong) & 0 \\
Wyler (1969) & Symmetric space & 137.036 & 0 \\
Standard Model & Fitted & 137.036 & 1 \\
RS (3-term) & Additive, $D = 3$ & 137.035 ($8.1\;\mathrm{ppm}$ off) & 0 \\
\bottomrule
\end{tabular}
\end{center}

\begin{remark}
Wyler's formula~\cite{Wyler1969} gives the correct numerical
value from symmetric-space volumes but lacks a physical
derivation.  RS provides both the correct value (at
$8.1\;\mathrm{ppm}$) and a physical mechanism: the geometric
seed minus the $J$-cost gap minus the curvature seam.
\end{remark}

%% ============================================================
\section{Predictions and Falsifiers}
\label{sec:predictions}
%% ============================================================

\begin{enumerate}
  \item \textbf{$\alpha^{-1} = 137.0349$ (Thomson limit)}:
  Zero-parameter prediction.  The $8.1\;\mathrm{ppm}$ residual
  from CODATA is attributed to higher-order corrections not
  included in the three-term formula.

  \item \textbf{Closed-form $w_8$}: The gap weight
  $w_8 = (246 + 145\sqrt{2} - 102\sqrt{5} - 65\sqrt{10})/7$
  is an algebraic number, not a fitted parameter.  Any independent
  derivation yielding a different $w_8$ would falsify the 8-tick
  epoch structure.

  \item \textbf{Running}: $\alpha$ runs from $\sim 1/137$ at
  $Q^2 = 0$ to $\sim 1/128$ at $M_Z$, as determined by QED
  vacuum polarization with RS-derived lepton masses.

  \item \textbf{Falsifier}: A future measurement showing
  $\alpha^{-1}(0) < 137.03$ or $\alpha^{-1}(0) > 137.04$ would
  rule out the three-term formula.
\end{enumerate}

%% ============================================================
\section{Conclusion}
\label{sec:conclusion}
%% ============================================================

The fine structure constant is derived from $D = 3$ cube geometry
via three additive terms: a geometric seed ($44\pi$), a gap
subtraction ($w_8\ln\varphi$), and a curvature correction
($103/(102\pi^5)$).  All integers---$11$, $6$, $17$, $103$---are
determined by the $D = 3$ cube, and the gap weight $w_8$ is a
parameter-free algebraic number.  The mystery of $137$ is
resolved: it is the geometric seed of passive photon channels,
reduced by the $J$-cost of epoch traversal and the topological
seam.

%% ============================================================
\begin{thebibliography}{99}

\bibitem{RS_Axioms}
J.~Washburn, ``The Algebra of Reality: A Recognition Science Derivation
of Physical Law,'' \emph{Axioms} \textbf{15}(2), 90 (2026).
\texttt{doi:10.3390/axioms15020090}

\bibitem{Morel2020}
L.~Morel, Z.~Yao, P.~Clad\'e, and S.~Guellati-Kh\'elifa,
``Determination of the Fine-Structure Constant with an Accuracy
of 81 Parts per Trillion,'' Nature \textbf{588}, 61 (2020).

\bibitem{Wyler1969}
A.~Wyler, ``L'espace sym\'etrique du groupe des \'equations de
Maxwell,'' C.~R.~Acad.~Sci.~Paris \textbf{269}, 743 (1969).

\bibitem{CODATA2024}
E.~Tiesinga, P.~J.~Mohr, D.~B.~Newell, and B.~N.~Taylor,
``CODATA Recommended Values of the Fundamental Physical
Constants: 2022,'' Rev.\ Mod.\ Phys.\ \textbf{96}, 025004
(2024).

\bibitem{Aczel1966}
J.~Acz\'el, \emph{Lectures on Functional Equations and Their
Applications} (Academic Press, 1966).

\end{thebibliography}

\end{document}
